problem:  
Let \((M^n, g, f)\) be a complete noncompact **steady gradient Ricci soliton** satisfying \(\mathrm{Ric}(g) = \nabla^2 f\), and assume the scalar curvature \(R > 0\) everywhere and \(|\mathrm{Rm}|\) is bounded.  
Suppose further that the potential function \(f\) has linear growth at infinity, i.e.  
\[
f(x) \sim a\, r(x) + b, \quad \text{for some constant } a > 0,
\]
where \(r(x)\) is the distance from a fixed origin.  

It is known that in all explicit examples (e.g. Bryant soliton), the **asymptotic cone** of \(M\) defined by the rescaling \( (M, r^{-2} g) \) as \(r \to \infty\) degenerates to a **half-line** (a 1-dimensional cone).  

**Question:**  
Does there exist a complete, nonflat, steady gradient Ricci soliton for which the asymptotic cone has dimension strictly greater than one—i.e. whose tangent cone at infinity is nontrivial and has a positive-dimensional base \( (N^{k}, h) \), \( k \ge 1 \)?  
Equivalently, can there exist steady solitons whose asymptotic geometry is **genuinely conical** (not collapsed to a ray), or is asymptotic 1-dimensional collapse an unavoidable universal feature of all nonflat steady solitons with bounded curvature?  

If such a conical-at-infinity steady soliton exists, describe the geometric constraints that its cross-section \(N^{k}\) must satisfy. If not, seek a rigidity principle establishing that every nonflat steady soliton must asymptotically collapse to a ray under appropriate rescaling.

Why is it a "good" problem:  
This problem probes the **asymptotic geometric dimension** of steady Ricci solitons, asking whether the known examples (all collapsing to 1-dimensional cones) reveal a universal collapse phenomenon or merely a special case. It reaches into the heart of the **global geometry of noncompact Ricci solitons**, testing whether the soliton equation enforces asymptotic dimensional reduction or allows new, higher-dimensional conical geometries.  

It is "good" because it directly bridges **geometric analysis**, **asymptotic geometry**, and **metric measure collapse theory**, raising a simple yet profound question: *can steady Ricci solitons open asymptotically like higher-dimensional cones, or must they always taper into one-dimensional geometry?*  
A resolution would illuminate the fundamental large-scale structure of steady solitons and deepen understanding of singularity models in the Ricci flow, with deep consequences for both geometric PDE theory and differential geometry.